\documentclass[a4paper,12pt]{article}

\usepackage[T2A]{fontenc}
\usepackage[utf8]{inputenc}
\usepackage[russian]{babel}

\usepackage{amsmath, amssymb, amsthm, mathtools}
\usepackage{helvet}
\renewcommand{\familydefault}{\sfdefault}
\usepackage{icomma}
\usepackage{geometry}
\usepackage{microtype}
\usepackage{tikz}
\usepackage{tkz-euclide}
\usepackage{pgfplots}
\pgfplotsset{compat=1.18}
\usepackage{tabularx, booktabs, longtable}
\usepackage{enumitem}
\usepackage{hyperref}
\usepackage{parskip}
\usepackage{fancyhdr}
\usepackage{setspace}
\usepackage{xcolor}

\definecolor{accent}{HTML}{008080}
\definecolor{darkaccent}{HTML}{004D4D}
\definecolor{textgray}{HTML}{333333}

\geometry{left=2.5cm, right=2.5cm, top=2.5cm, bottom=2.5cm}
\onehalfspacing

\begin{document}

% Титульный лист
\begin{titlepage}
\centering
\vspace*{2cm}
{\Large\bfseries\color{darkaccent} Чек-лист}\par
\vspace{0.5cm}
{{large Тема: Показательные уравнения и параметры}}\par
\vspace{2cm}
{\large Лёвин Артём Александрович}\par
{\large эксклюзивно для Мика}\par
\vspace{2cm}
{\today}\par
\vfill
\begin{tikzpicture}
\draw[color=accent, line width=1.5pt] 
(0,0) .. controls (2,1.5) and (4,-0.5) .. (6,0.5) .. controls (8,1.5) and (10,-0.5) .. (12,0);
\end{tikzpicture}
\vspace{1cm}
\end{titlepage}

\section*{Чек-лист: Показательные уравнения и параметры}

\subsection*{Основные понятия}
\item \textbf{Показательное уравнение}: уравнение вида $a^{f(x)} = b^{g(x)}$, где $a,b > 0$, $a,b \neq 1$.
\item \textbf{Метод замены}: введение новой переменной $t = a^{f(x)}$.
\item \textbf{Метод логарифмирования}: применение логарифма к обеим частям уравнения.
\item \textbf{ОДЗ для показательной функции}: $a^{f(x)} > 0$ при любых $x$.

\subsection*{Методы решения}
\begin{enumerate}[leftmargin=*]
\item Приводить к одному основанию степени.
\item Выносить общий множитель за скобки.
\item Делить на положительную степень (не равную нулю).
\item Для уравнений с параметром использовать графический метод или анализ условий.
\item Проверять полученные корни на соответствие ОДЗ.
\end{enumerate}

\subsection*{Типичные ошибки}
\begin{itemize}[leftmargin=*]
\item[$\times$] Делят на выражение, которое может быть равно нулю.
\item[$\times$] Забывают проверить принадлежность корней области определения.
\item[$\times$] Неправильно применяют свойства логарифмов.
\item[$\times$] При замене переменной забывают обратную замену.
\item[$\times$] Не учитывают монотонность показательной функции.
\end{itemize}

\newpage
\section*{Тренировочный блок}

\textbf{Уровень 1 (базовый)}

1. Решите уравнение: $2^{3x} = 8$.

2. Решите: $3^{x+1} = 27$.

3. Найдите область определения: $y = \sqrt{2^x - 1}$.

\medskip
\textbf{Уровень 2 (средний)}

4. Решите: $2^{2x} - 5\cdot 2^x + 4 = 0$.

5. Решите уравнение: $3^{x} \cdot 5^{x} = 15$.

6. Найдите число корней уравнения $2^x = ax$ при $a = 1$.

\medskip
\textbf{Уровень 3 (выше среднего)}

7. Решите: $4^x - 3\cdot 2^{x+1} + 8 = 0$.

8. При каких значениях параметра $a$ уравнение $2^x = a$ имеет единственное решение?

9. Решите неравенство: $\log_2(2^x - 1) < 3$.

\medskip
\textbf{Уровень 4 (сложный)}

10. Найдите все значения параметра $a$, при которых уравнение $4^x - a\cdot 2^x + a^2 - 1 = 0$ имеет ровно один корень.


\end{document}
